% Primary slide source. Compile through scripts/build_introduction_pdf.py.
% Cairo extracts and exported PDFs remain local-only; see NOTICE.md.
\documentclass[aspectratio=169,11pt]{beamer}
\usepackage{fontspec}
\usepackage{unicode-math}
\setsansfont{Arial}
\setmonofont{Menlo}
\setmathfont{latinmodern-math.otf}
\usepackage{booktabs,tabularx,pgfplots}
\pgfplotsset{compat=1.18}
\definecolor{ink}{HTML}{172D43}
\definecolor{accent}{HTML}{176B73}
\definecolor{muted}{HTML}{536270}
\definecolor{orange}{HTML}{B54C22}
\definecolor{pale}{HTML}{F4D4C3}
\definecolor{blue}{HTML}{4363A8}
\setbeamercolor{normal text}{fg=ink,bg=white}
\setbeamercolor{frametitle}{fg=ink}
\setbeamercolor{structure}{fg=accent}
\setbeamerfont{frametitle}{size=\Large,series=\bfseries}
\setbeamertemplate{navigation symbols}{}
\setbeamersize{text margin left=8mm,text margin right=8mm}
\setbeamertemplate{frametitle}{\vspace{3mm}\insertframetitle\par}
\setbeamertemplate{footline}{\hspace{8mm}\color{muted}\fontsize{6}{7}\selectfont
ECBS5209 · Data Visualization and Storytelling · Ulrich Wohak\hfill\insertframenumber\hspace{8mm}\vspace{3mm}}
\hypersetup{colorlinks=true,urlcolor=accent,pdftitle={Data Visualization and Storytelling},pdfauthor={Ulrich Wohak}}
\newcommand{\point}[1]{\par\vspace{4mm}#1\par}
\newcommand{\takeaway}[1]{\par\vfill\textcolor{accent}{\textbf{#1}}\par\vspace{2mm}}
\newcommand{\smallnote}[1]{\par\vspace{2mm}{\scriptsize\color{muted}#1}\par}
\begin{document}

\begin{frame}
\vspace{7mm}
{\LARGE “The purpose of visualization is\par}
\vspace{3mm}
{\huge\bfseries\color{accent} insight, not pictures.”\par}
\vfill
{\Large\bfseries Data Visualization and Storytelling\par}
\vspace{4mm}
ECBS5209 · CEU · Fall 2026\par
Ulrich Wohak · 4 November 2026\par\vspace{4mm}
\end{frame}

\begin{frame}{Welcome to the course}
\point{Ulrich Wohak, PhD · \href{mailto:wohaku@ceu.edu}{wohaku@ceu.edu}}
\point{MSc Business Analytics · 1 US credit / 2 ECTS}
\point{Six Wednesday sessions · 17:40–19:20}
\point{Bring your laptop and your existing Python and pandas skills.}
\end{frame}

\begin{frame}{What you should be able to do}
\point{Frame a useful question for a particular audience.}
\point{Choose and critique charts that make comparison easier.}
\point{Tell a clear story while retaining uncertainty and context.}
\point{Build and verify visualisations in Python and Power BI.}
\takeaway{Design judgment matters as much as technical execution.}
\end{frame}

\begin{frame}{How we will work in class}
\point{Discuss a question and inspect the evidence.}
\point{Build Python examples, then Power BI charts and dashboards.}
\point{Inspect a chart and decide what you would change.}
\point{Explain your choices and check the result.}
\takeaway{The code may run perfectly while the chart communicates badly.}
\end{frame}

\begin{frame}{Our six sessions}
\fontsize{8.5}{10.2}\selectfont
\hyphenpenalty=10000
\renewcommand{\arraystretch}{1.18}
\begin{tabularx}{\textwidth}{@{}p{13mm}X@{}}
\textbf{Date} & \textbf{Focus}\\\midrule
4 Nov & \textbf{Python chart tour:} line, bar, scatter, histogram, box, violin and strip plots; labels, notes and export.\\[1.4mm]
11 Nov & \textbf{Comparisons and uncertainty:} histograms, box plots and confidence-interval plots, including a supplied event-study example.\\[1.4mm]
18 Nov & \textbf{Mapping in Python:} geographical joins, choropleths, proportional symbols, projections and meaningful denominators.\\[1.4mm]
25 Nov & \textbf{Power BI foundations:} interface, importing data, field types, aggregation and basic charts.\\[1.4mm]
2 Dec & \textbf{Power BI comparisons and filtering:} tables, matrices, cards/KPIs, filters and slicers.\\[1.4mm]
9 Dec & \textbf{Power BI dashboard building:} layout, interactions, storytelling and checking filtered results.\\
\end{tabularx}
\smallnote{Wednesdays, 17:40–19:20 · 2026}
\end{frame}

\begin{frame}{Assessment}
\vspace{6mm}
{\Huge\bfseries\color{accent}25\%}\hspace{8mm}{\large Three individual figures with notes}\par
\smallnote{Equal weight within the combined 25\%}
\vspace{9mm}
{\Huge\bfseries\color{accent}75\%}\hspace{8mm}{\large Final Python data story}\par
\vfill
In-class exercises: practice and feedback\par\vspace{3mm}
Minimum overall mark to pass: 50\%
\smallnote{Submission dates and routes: to be announced}
\end{frame}

\begin{frame}{Three individual figures · 25\% in total}
\point{Three submissions with equal weight.}
\point{For each, I supply the data and the question.}
\point{You design a figure in Python that answers it.}
\point{Submit only the figure with its accompanying notes.}
\takeaway{Data, questions and deadlines will be announced later.}
\end{frame}

\begin{frame}{The final Python data story · 75\%}
\point{One coherent question and a specific audience.}
\point{A sequence of 3–5 figures that work together.}
\point{Reproducible code, data provenance and a written narrative.}
\point{Work individually or in a team of up to three.}
\takeaway{A strong static story is enough. Interactivity is optional.}
\end{frame}

\begin{frame}{Where to find the material}
\href{https://github.com/ulrichwohak/Python-Dataviz}{\large github.com/ulrichwohak/Python-Dataviz}
\point{One teaching notebook and one exercise notebook per lecture}
\point{Python: pandas, matplotlib and seaborn · environment: \texttt{uv}}
\smallnote{No required textbook purchase. Readings will be identified in class.}
\point{Explore datasets: \href{https://github.com/toddrjones/datasets}{\underline{github.com/toddrjones/datasets}}}
\end{frame}

\begin{frame}{What are we trying to make small?}
\vspace{4mm}
\[
 (\widehat{\beta}_0,\widehat{\beta}_1)
 =\underset{\beta_0,\beta_1}{\operatorname{arg\,min}}
 \sum_{i=1}^{n}\left[y_i-(\beta_0+\beta_1 x_i)\right]^2
\]
\vspace{6mm}
\[
 \widehat{\symbf{\beta}}
 =\underset{\symbf{\beta}}{\operatorname{arg\,min}}
 (\symbf{y}-\symbf{X}\symbf{\beta})^{\mathsf T}
 (\symbf{y}-\symbf{X}\symbf{\beta})
\]
\takeaway{The equation gives precision. Can a picture give intuition?}
\end{frame}

\begin{frame}[t]{A relationship with variation}
\vspace{2mm}
\centering
\input{ols-scatter.tex}\par
\smallnote{100 synthetic observations}
\end{frame}

\begin{frame}[t]{The OLS fitted line}
\vspace{2mm}
\centering
\input{ols-fitted.tex}\par
\smallnote{100 synthetic observations. The line minimises the sum of squared vertical residuals.}
\takeaway{“We understand \emph{because} we see.” — Alberto Cairo, p. xv}
\end{frame}

\begin{frame}{What would you conclude from this chart?}
\begin{columns}[c,onlytextwidth]
\column{.68\textwidth}\input{simpson-pooled.tex}
\column{.29\textwidth}
\textcolor{accent}{\textbf{Across all hotels:\\price rises with distance.}}
\vspace{10mm}\par What information might be missing?
\end{columns}
\smallnote{Synthetic illustration · 24 hotels · three fictional countries · association, not causation}
\end{frame}

\begin{frame}{The same points tell a different story}
\begin{columns}[c,onlytextwidth]
\column{.68\textwidth}\input{simpson-grouped.tex}
\column{.29\textwidth}
\textcolor{accent}{\textbf{Within each country:\\price falls with distance.}}
\vspace{6mm}\par
\textcolor{accent}{$\bullet$ Country A}\par\vspace{2mm}
\textcolor{orange}{$\blacksquare$ Country B}\par\vspace{2mm}
\textcolor{blue}{Country C}
\end{columns}
\smallnote{Synthetic illustration · same points and scales · association, not causation}
\end{frame}

\begin{frame}{What is this figure trying to communicate?}
\begin{columns}[c,onlytextwidth]
\column{.59\textwidth}\centering
\includegraphics[height=57mm,width=\linewidth,keepaspectratio]{cairo-figure3-full.png}
\column{.37\textwidth}
\textcolor{accent}{\large\textbf{Describe the takeaway in one sentence.}}
\vspace{9mm}\par Would another reader choose the same message?
\end{columns}
\smallnote{Cairo / Época · \emph{The Functional Art}, Figure 3, p. xix}
\end{frame}

\begin{frame}{Which details serve the intended purpose?}
\begin{columns}[c,onlytextwidth]
\column{.70\textwidth}\centering
\includegraphics[height=57mm,width=\linewidth,keepaspectratio]{cairo-figure3-chart.png}
\column{.27\textwidth}
\textcolor{accent}{\textbf{A focused finding?}}\par\vspace{7mm}
\textcolor{accent}{\textbf{Historical exploration?}}\par\vspace{7mm}
Keep, remove, or separate?
\end{columns}
\smallnote{Cairo / Época · Figure 3, p. xix · same figure, enlarged chart detail}
\end{frame}

\begin{frame}{Each figure needs a purpose and context}
\point{Who is reading, and what are they trying to understand?}
\point{What should this figure reveal, support or let them explore?}
\point{What would become misleading if we left it out?}
\point{Keep necessary units, sample, uncertainty and source notes.}
\takeaway{Simplify the presentation while preserving the evidence.}
\end{frame}

\begin{frame}{A quick Reddit chart critique}
\href{https://www.reddit.com/r/dataisugly/}{\large\underline{reddit.com/r/dataisugly}}\par\vspace{4mm}
\href{https://www.reddit.com/r/dataisbeautiful/}{\large\underline{reddit.com/r/dataisbeautiful}}\par\vspace{5mm}
\textbf{Choose a chart online or from the following examples.}
\point{What do you learn from it?}
\point{What is clear, and what could mislead the reader?}
\point{What one change would you make?}
\end{frame}

\begin{frame}{Prices and wages}
\begin{columns}[c,onlytextwidth]
\column{.68\textwidth}\centering
\includegraphics[height=61mm,width=\linewidth,keepaspectratio]{prices-and-wages.png}
\column{.29\textwidth}
\textcolor{accent}{\textbf{What comparison does this picture invite?}}
\vspace{7mm}\par What would you check before accepting it?
\end{columns}
\smallnote{Instructor-supplied photo showing CBS News and a BLS credit. Figures not independently checked.}
\end{frame}

\begin{frame}{Reading a time series}
\begin{columns}[c,onlytextwidth]
\column{.62\textwidth}\centering
\includegraphics[height=61mm,width=\linewidth,keepaspectratio]{florida-firearm-murders.png}
\column{.35\textwidth}
\textcolor{accent}{\textbf{What happened after 2005?}}
\vspace{7mm}\par Does your first impression agree with the scale?
\end{columns}
\smallnote{Instructor-supplied image credited to Reuters / C. Chan, 16 February 2014.}
\end{frame}

\begin{frame}{A pool and three skyscrapers}
\begin{columns}[c,onlytextwidth]
\column{.66\textwidth}\centering
\includegraphics[height=61mm,width=\linewidth,keepaspectratio]{pool-and-skyscrapers.png}
\column{.31\textwidth}
\textcolor{accent}{\textbf{What does “bigger” mean here?}}
\vspace{7mm}\par Do the labels and picture support the same comparison?
\end{columns}
\smallnote{Instructor-supplied social-media screenshot. Original post and measurements not verified.}
\end{frame}

\begin{frame}{Reading a county map}
\begin{columns}[c,onlytextwidth]
\column{.66\textwidth}\centering
\includegraphics[height=61mm,width=\linewidth,keepaspectratio]{county-suicide-rates.png}
\column{.31\textwidth}
\textcolor{accent}{\textbf{What do the colours tell you?}}
\vspace{7mm}\par What else would you need to interpret this map?
\end{columns}
\smallnote{Instructor-supplied screenshot labelled amazingmap. Underlying source and data period not shown.}
\end{frame}

\begin{frame}{Today: a chart tour in Python}
\textcolor{accent}{\textbf{US birth counts: what changes across weekdays?}}\par\vspace{2mm}
\point{Line, bar and scatter plots for concrete questions.}
\point{Histograms, box plots, violins and individual data points.}
\point{Add labels and figure notes, then export a figure.}
\point{Final 20 minutes: redesign the supplied chart and check the values.}
\takeaway{Keep the question visible while we write the code.}
\end{frame}

\begin{frame}[fragile]{Backup · setup and Power BI access}
\textcolor{accent}{\textbf{Python: from the repository root}}
\begin{verbatim}
uv sync
uv run python scripts/fetch_data.py
uv run jupyter lab
\end{verbatim}
\vspace{2mm}
\smallnote{First data download: about 330 MB. Later runs use verified local caches.}
\vspace{2mm}
Before Session 4: Power BI access will be confirmed.\par\vspace{3mm}
Windows / CEU access route: to be announced.\par\vspace{3mm}
No purchase is required for this course.
\end{frame}
\end{document}
